\documentclass[11pt]{article}
\usepackage[margin=0.9in]{geometry}
\usepackage{amsmath,array,booktabs,enumitem,tikz}
\usepackage[version=4]{mhchem}
\usetikzlibrary{arrows.meta}
\setlength{\parindent}{0pt}\setlength{\parskip}{6pt}
\renewcommand{\arraystretch}{1.5}
\pagestyle{plain}
% orbital box: 0 empty, 1 up, 2 paired
\newcommand{\ob}[1]{\tikz[baseline=-0.7ex]{\draw (0,-0.28) rectangle (0.5,0.28);%
\ifcase#1\or\draw[-{Stealth[length=3pt]}](0.25,-0.2)--(0.25,0.2);%
\or\draw[-{Stealth[length=3pt]}](0.17,-0.2)--(0.17,0.2);\draw[-{Stealth[length=3pt]}](0.33,0.2)--(0.33,-0.2);\fi}}
\newcommand{\sub}[2]{\begin{tabular}[b]{@{}c@{}}#2\\[-6pt]\small #1\end{tabular}\hspace{10pt}}
\newcommand{\pb}[3]{\ob{#1}\hspace{-0.4pt}\ob{#2}\hspace{-0.4pt}\ob{#3}}
\newcommand{\ans}[1][1]{\par\foreach \i in {1,...,#1}{\vspace{0.8cm}\hrule}\vspace{4pt}}
\newcommand{\nm}{Name: \rule{7cm}{0.4pt}\par}

\begin{document}
{\Large\bfseries Review Practice: Electrons, Ions and Groups}\par\nm
\section*{Electrons}
\textbf{1.} Write the electron configuration.

a.\ sodium (Na), 11 electrons: \quad $1s^2\,2s^2\,2p^6$ \rule{3cm}{0.4pt}\par\vspace{6pt}
b.\ magnesium (Mg), 12 electrons: \rule{7cm}{0.4pt}\par\vspace{6pt}
c.\ chlorine (Cl), 17 electrons: \rule{7.3cm}{0.4pt}

\textbf{2.} Fill in the orbital boxes for oxygen (8 electrons).\par\vspace{6pt}
\sub{1s}{\ob0}\sub{2s}{\ob0}\sub{2p}{\pb000}

\textbf{3.} Complete the table.

\begin{tabular}{|c|p{5cm}|c|c|c|}\hline
Element & Configuration & Energy levels & Core $e^-$ & Valence $e^-$\\\hline
S & $1s^2\,2s^2\,2p^6\,3s^2\,3p^4$ & 3 & 10 & 6\\\hline
Na & & & & \\\hline
Cl & & & & \\\hline
K & $1s^2\,2s^2\,2p^6\,3s^2\,3p^6\,4s^1$ & & & \\\hline
\end{tabular}

\section*{Cations and anions}
\textbf{4.} Complete the table.

\begin{tabular}{|c|c|c|c|}\hline
Ion & Protons & Electrons & Cation or anion?\\\hline
\ce{Mg^2+} & 12 & 10 & cation\\\hline
\ce{Na+} & & & \\\hline
\ce{O^2-} & & & \\\hline
\ce{Br-} & & & \\\hline
\ce{Al^3+} & & & \\\hline
\end{tabular}

\textbf{5.} A calcium atom loses 2 electrons. Write the ion's symbol and say whether it is a cation or an anion.\ans

\section*{Groups}
\textbf{6.} Complete the table.

\begin{tabular}{|c|c|p{4cm}|c|c|}\hline
Element & Group & Group name & Valence $e^-$ & Ion\\\hline
Li & 1 & Alkali metals & 1 & \ce{Li+}\\\hline
Ca & & & & \\\hline
F & & & & \\\hline
Ne & & & & \\\hline
\end{tabular}

\textbf{7.} Use this list: K, F, Mg, Ar, Ca, Br. Which element or elements:

\begin{tabular}{@{}ll@{}}
a.\ are halogens? \rule{3cm}{0.4pt} & b.\ is a noble gas? \rule{3cm}{0.4pt}\\
c.\ are alkaline earth metals? \rule{3cm}{0.4pt} & d.\ is an alkali metal? \rule{3cm}{0.4pt}
\end{tabular}

\textbf{8.} Which is more reactive, potassium or calcium? Explain using valence electrons.\ans[2]

\textbf{9.} Why does argon almost never react?\ans
\end{document}
